\documentclass[11pt]{article}
\usepackage{explorer}
\newcommand{\ExplorerLanguage}{EN}
\newcommand{\ExplorerNumber}{05}
\newcommand{\ExplorerKind}{Lesson and method}
\begin{document}
\ExplorerTitle{High school · counting}{Paths, combinations and Pascal}
\ExplorerSection{Before you begin}
Products, factorials and unordered choices.\par
\ExplorerSection{Ideas and method}
A path from (0,0) to (a,b), using only rightward or upward steps, has \(a + b\) steps. Choosing the positions of the a horizontal steps gives \(C(a+b,a)\) paths.\par
For \(0 \leq  k \leq  n, C(n,k) = n\)!/[k!\((n-k)\)!]. The order of the chosen objects does not matter. Symmetry \(C(n,k) = C(n,n-k)\) corresponds to choosing included or excluded objects.\par
Pascal’s identity \(C(n,k) = C(n-1,k-1) + C(n-1,k)\) splits choices by whether one fixed object is included. Paths can also be classified by their last step.\par
\ExplorerSection{Worked example}
From (0,0) to (3,2), there are \(C(5,3) = 10\) paths. To pass through (1,1), choose one of two paths to that point and one of three paths onward: \(2 \times  3 = 6\) paths.\par
\begin{center}\begin{tikzpicture}[scale=.85]\draw[gray!45](0,0)grid(3,2);\draw[mint,very thick,->](0,0)--(1,0)--(1,1)--(2,1)--(3,1)--(3,2);\fill[violet](1,1)circle(3pt)node[above left]{$(1,1)$};\node[below]at(0,0){$(0,0)$};\node[above]at(3,2){$(3,2)$};\end{tikzpicture}\end{center}
\ExplorerSection{Explore}
Draw a small grid and write the number of arriving paths at each vertex. Add the values of its left and lower neighbours.\par
\end{document}
